
\begin{obereflection}
\textbf{Latihan 5.1}
1. Jika plainteks adalah \texttt{10110010} dan kunci adalah \texttt{11001010}, hitunglah cipherteksnya.
2. Mengapa \textit{stream cipher} lebih cocok digunakan untuk komunikasi suara (VoIP) dibandingkan \textit{block cipher}?
3. Jelaskan mengapa penggunaan kunci yang sama untuk dua pesan berbeda dalam \textit{stream cipher} (disebut \textit{key reuse}) sangat berbahaya.
4. Nyatakan kalimat \texttt{Kripto} dalam bentuk biner ASCII, lalu tuliskan pula bentuk heksadesimalnya.
5. Sebuah LFSR 4 bit bermuatan awal \texttt{1010} memakai umpan balik $f=b_3\oplus b_0$. Tentukan isi register pada empat detak pertama beserta bit keluarannya.
6. Mengapa LFSR tunggal tidak aman dipakai sebagai \textit{keystream generator}, padahal periode keluarannya dapat mencapai $2^n-1$ bit?
7. Jelaskan perbedaan peran \textit{key}, \textit{nonce}, dan \textit{counter} pada Salsa20. Manakah yang harus dirahasiakan, dan mengapa?
8. Bandingkan Salsa20 dan ChaCha20 dari segi bentuk \textit{quarter round} serta akibatnya pada penyebaran bit.
\end{obereflection}

\begin{obereflection}
\textbf{Latihan 5.2: Hitungan dan Analisis}
1. Sebuah pesan 12 bit \texttt{011011010110} dienkripsi dengan LFSR 4 bit bermuatan awal \texttt{1111} dan umpan balik $f=b_3\oplus b_0$. Tuliskan keystream 12 bit yang dipakai, lalu tentukan cipherteksnya.
2. Pada RC4 versi mainan dengan $N=8$ dan kunci \texttt{B}, tentukan larik $S$ setelah KSA, lalu hitung byte keystream pertama.
3. Frame GSM berukuran 228 bit dikirim setiap 4,6 milidetik. Hitunglah laju bit yang harus dihasilkan A5/1 untuk mengikuti aliran itu. Mengapa \textit{warm-up} 100 putaran tidak mengganggu laju tersebut?
4. Trivium memakai tiga NLFSR berukuran 93, 84, dan 111 bit dengan total keadaan 288 bit. Jelaskan mengapa penambahan gerbang AND membuat serangan berbasis persamaan linier tidak lagi bekerja.
5. Bandingkan jumlah \textit{warm-up clock} A5/1 dan Trivium relatif terhadap jumlah bit yang dihasilkan per siklus. Apa alasan perbedaan itu?
6. Seorang penyerang menangkap tiga cipherteks TLS yang seluruhnya memakai \textit{nonce} ChaCha20 yang sama tetapi kunci berbeda. Apakah penyerang memperoleh keuntungan? Jelaskan.
7. Seorang penyerang menangkap dua cipherteks $c_1$ dan $c_2$ yang dibangkitkan dengan keystream yang sama, serta berhasil menebak bahwa plainteks pertama $p_1$ dimulai dengan kata \texttt{LAPORAN}. Tunjukkan bagaimana $p_2$ dapat dipulihkan tanpa mengetahui kuncinya.
\end{obereflection}
